\documentclass[11pt,a4paper]{report}
\usepackage[T1]{fontenc}
\usepackage[utf8]{inputenc}
\usepackage{lmodern}
\usepackage{amsmath,graphicx,array,xcolor}
\IfFileExists{french.ldf}{\usepackage[french]{babel}}{%
  \renewcommand{\contentsname}{Table des matières}
  \renewcommand{\chaptername}{Chapitre}
  \renewcommand{\appendixname}{Annexe}
  \renewcommand{\bibname}{Références bibliographiques}
  \renewcommand{\figurename}{Figure}
  \renewcommand{\tablename}{Tableau}}
\definecolor{accent}{HTML}{173C59}
\usepackage[colorlinks=true,linkcolor=accent,citecolor=accent,urlcolor=accent]{hyperref}
\setlength{\textwidth}{157mm}
\setlength{\oddsidemargin}{1.1mm}
\setlength{\evensidemargin}{1.1mm}
\setlength{\textheight}{230mm}
\setlength{\topmargin}{-10mm}
\setlength{\parskip}{0.25em}
\setlength{\parindent}{1em}
\setcounter{tocdepth}{1}
\makeatletter
\renewcommand{\@makechapterhead}[1]{%
  \vspace*{20pt}%
  {\parindent\z@\raggedright\normalfont
   \ifnum\c@secnumdepth>\m@ne
     \large\bfseries\color{accent}\@chapapp\space\thechapter\par\nobreak
     \vskip 8pt
   \fi
   \interlinepenalty\@M\huge\bfseries #1\par\nobreak\vskip 20pt}}
\renewcommand{\@makeschapterhead}[1]{%
  \vspace*{20pt}%
  {\parindent\z@\raggedright\normalfont\interlinepenalty\@M
    \huge\bfseries #1\par\nobreak\vskip 20pt}}
\makeatother

\graphicspath{{figures/}}
\newcommand{\nomEtudiant}{[Nom et prénom à compléter]}
\newcommand{\nomEncadrant}{[Nom et prénom à compléter]}
\newcommand{\anneeUniversitaire}{[Année universitaire à compléter]}
\newcommand{\sourcetable}[1]{\par\smallskip{\footnotesize\emph{Source :} #1\par}}
\newcommand{\figureprojet}[3]{%
  \begin{figure}[!htbp]
  \centering\includegraphics[width=\linewidth,height=.70\textheight,keepaspectratio]{#1}
  \caption{#2}\label{#3}
  \end{figure}}
\hypersetup{pdftitle={Inflation \& Consumer Behavior Analysis - Morocco}}

\begin{document}
\hypersetup{pageanchor=false}
\begin{titlepage}
\centering
{\Large Institut National de Statistique et d'Économie Appliquée\par}
\medskip
{\large INSEA\par}
\vspace{2.4cm}
{\large Rapport de projet académique individuel\par}
\medskip
{\normalsize Première année -- Économie Appliquée, Statistique et Big Data\par}
\vspace{1.25cm}
{\color{accent}\rule{0.82\textwidth}{0.9pt}\par}
\vspace{0.65cm}
{\huge\bfseries Inflation \& Consumer Behavior\\[0.22em] Analysis -- Morocco\par}
\vspace{0.5cm}
{\large Prix à la consommation et exposition des budgets des ménages marocains\par}
\vspace{0.65cm}
{\color{accent}\rule{0.82\textwidth}{0.9pt}\par}
\vfill
\begin{tabular}{@{}ll@{}}
Étudiant & \nomEtudiant \\[0.7em]
Année universitaire & \anneeUniversitaire
\end{tabular}
\vfill
\end{titlepage}

\hypersetup{pageanchor=true}
\pagenumbering{roman}
\chapter*{Résumé}
\addcontentsline{toc}{chapter}{Résumé}
Ce rapport étudie l'évolution des prix au Maroc et l'exposition de différents profils budgétaires à ces variations. Il combine 118 observations mensuelles de l'indice des prix à la consommation (IPC) du Haut-Commissariat au Plan, de janvier 2017 à octobre 2026, et 18 profils agrégés de dépenses de l'Enquête nationale sur le niveau de vie des ménages 2022--2023. L'IPC général atteint un pic de variation mensuelle de 1,79\,\% en mars 2022 et un pic en glissement annuel de 10,08\,\% en février 2023. Pour l'ensemble national, le coût simulé à quantités fixes augmente entre septembre 2022 et octobre 2026 de 9,22\,\% dans le sous-panier prudent et de 7,85\,\% dans le sous-panier élargi. Ces périmètres couvrent respectivement 41,4\,\% et 93,2\,\% du budget publié. Les simulations hypothétiques montrent comment la couverture budgétaire et la part des aliments déterminent le supplément à financer à quantités constantes. Pour l'ensemble national, sous l'hypothèse d'une hausse nominale de 5\,\%, le nombre d'unités accessibles en octobre 2026 diminue par rapport à septembre 2022 de 3,86\,\% dans le prudent et de 2,64\,\% dans l'élargi. Les profils ruraux sont soumis aux mêmes séries de prix à collecte urbaine que les autres profils. L'analyse estime une exposition conditionnelle aux prix. La mesure des changements de consommation et du pouvoir d'achat effectif exigerait des données complémentaires.

\noindent\textbf{Mots-clés :} Maroc, IPC, ENNVM, budget des ménages, panier fixe, prix alimentaires.
\tableofcontents
\clearpage
\pagenumbering{arabic}
\chapter*{Introduction générale}
\addcontentsline{toc}{chapter}{Introduction générale}
Une hausse des prix modifie la dépense nécessaire pour acheter une quantité donnée de biens et services. Son effet sur un budget dépend de la composition des achats. Lorsque l'alimentation occupe une part importante des dépenses, une poussée des prix alimentaires peut peser davantage sur l'équilibre budgétaire. Les données marocaines permettent d'étudier ce mécanisme à deux niveaux. L'IPC renseigne sur les prix suivis par le Haut-Commissariat au Plan (HCP). L'Enquête nationale sur le niveau de vie des ménages (ENNVM) décrit la répartition des dépenses de profils agrégés \cite{hcpIPC,hcpENNVM}.

L'exposition budgétaire désigne ici la variation du coût d'un sous-panier de composition fixe lorsque les prix changent. L'IPC renseigne sur les variations de prix et l'ENNVM 2022--2023 fournit les pondérations à un moment donné. Leur rapprochement répond à une question de simulation : quelle dépense supplémentaire faudrait-il engager pour conserver les mêmes quantités ? Observer le comportement des consommateurs demanderait de suivre leurs achats dans le temps pour mesurer, par exemple, des substitutions, des réductions de quantités ou des renoncements.

\section*{Problématique et objectifs}
Comment les différences de structure des dépenses modifient-elles l'exposition aux hausses de prix des profils nationaux, urbains et ruraux et des différents quintiles de dépenses au Maroc ? Le rapport décrit d'abord la dynamique de l'IPC général et de ses divisions. Il rapproche ensuite les postes de l'ENNVM des séries IPC disponibles, construit des indices de coût à poids fixes, puis distingue l'évolution aux prix historiques de deux chocs hypothétiques. La comparaison des deux périmètres permet d'apprécier dans quelle mesure les résultats dépendent de la couverture et des rapprochements de catégories.

\section*{Hypothèses de travail}
\begin{enumerate}
\item Une hausse uniforme de tous les prix retenus produit le même taux de variation pour chaque sous-panier. Le supplément imputable au budget total varie avec la part couverte.
\item Un choc alimentaire supplémentaire accroît particulièrement la dépense nécessaire pour les profils dont la part alimentaire initiale est élevée. Le résultat suppose des quantités fixes et un même budget initial dans l'illustration en dirhams.
\item L'élargissement du panier peut modifier le taux calculé et sa contribution au budget initial, puisque les nouveaux postes ont des trajectoires de prix propres et modifient les poids relatifs.
\end{enumerate}
Ces hypothèses décrivent des mécanismes comptables et économiques. Les données permettent de vérifier leur traduction dans les simulations.

\chapter{Cadre théorique}
\section{Niveau des prix, indice et inflation}
Un indice des prix à la consommation résume l'évolution des prix d'un ensemble de biens et services. La base 100 en 2017 fixe son échelle. Une valeur de 121,8 représente 121,8 points d'indice. Le taux d'inflation se calcule en comparant deux niveaux d'indice et s'exprime en pourcentage. Une inflation positive qui ralentit signifie que le niveau des prix reste supérieur à celui de la période de comparaison, avec un écart plus faible. Une variation négative indique une baisse pour l'horizon calculé.

Le HCP décrit un IPC de type Laspeyres en chaîne, avec une moyenne géométrique au premier niveau d'agrégation. Les pondérations reposent principalement sur les données de 2013--2014, complétées pour certains postes. La collecte porte sur 18 villes et le champ de référence est urbain \cite{hcpIPC,hcpNoteComplete}. Les sous-paniers du projet repondèrent des indices de divisions déjà agrégés. Leurs poids viennent de l'ENNVM 2022--2023.

\section{Contrainte budgétaire et choix de consommation}
Soit $p_k$ le prix unitaire d'un bien $k$, $q_k$ la quantité choisie et $Y$ le budget nominal consacré aux achats. L'ensemble des paniers accessibles satisfait
\[
  \sum_k p_kq_k \leq Y,\qquad q_k\geq 0.
\]
Une hausse uniforme des prix réduit les quantités accessibles à budget constant, avec des prix relatifs inchangés. Une hausse concentrée sur un poste modifie les prix relatifs et peut inciter à déplacer une partie des achats. La théorie distingue l'effet de substitution, lié aux prix relatifs, et l'effet de revenu, lié à la variation de la capacité d'achat réelle \cite{openstaxChoix}. À titre d'exemple, une hausse du prix d'un aliment peut inciter à choisir un produit moins cher. L'ampleur et la direction de la réponse totale dépendent des préférences, des possibilités de remplacement et du revenu.

\section{Coût du panier et pouvoir d'achat}
La dépense nécessaire pour acheter un panier donné $q^0$ est $E_t(q^0)=\sum_k p_{k,t}q_k^0$. Le rapport $E_t(q^0)/E_b(q^0)$ décrit un coût à quantités fixes. Un indice du coût de la vie au sens microéconomique demanderait de comparer les dépenses minimales nécessaires pour atteindre un même niveau de satisfaction aux deux dates, en autorisant la substitution. Le pouvoir d'achat d'un revenu demande par ailleurs un revenu nominal comparable et un déflateur adapté. Un IPC de base 100 permet d'exprimer un revenu nominal en prix de la période de base par $100Y_t/I_t$. Les fichiers exploités ici apportent des parts budgétaires agrégées et des indices de prix. La trajectoire du coût simulé permet d'isoler le canal des prix sous les hypothèses du modèle.

\chapter{Sources, préparation et traçabilité}
\section{Séries mensuelles d'IPC}
Les trois classeurs mensuels du HCP couvrent janvier 2017 à octobre 2026 \cite{hcpIndicateurs}. Les grandes divisions fournissent l'indice général, l'alimentation et les produits non alimentaires. Les deux groupes de divisions fournissent six séries chacun. Le fichier annuel de détail « Biens et services divers » éclaire la nomenclature. Il ne complète pas la série mensuelle.

Le script \texttt{clean\_ipc.py} lit les en-têtes à la ligne 24 et les observations à partir de la ligne 25. Il élimine les cellules de présentation et la note de source, transforme le mois \texttt{AAAA/MM} en date au premier du mois et trie la série par ordre croissant. Il contrôle les intitulés, les doublons de dates, la présence des 118 mois attendus et les valeurs numériques positives et finies. Une fusion sur la date associe les 118 lignes de chaque tableau. Les libellés HCP sont conservés, y compris la distinction entre l'agrégat « Alimentation » et sa division « Produits alimentaires et boissons non alcoolisées ». Le CSV final compte 118 lignes, 15 séries et une colonne date. Il ne comporte ni mois dupliqué, ni valeur manquante. Les 1\,770 indices ont été rapprochés des cellules Excel lors de la préparation.

\section{Enquête sur le niveau de vie des ménages}
L'ENNVM a été collectée du 15 mars 2022 au 14 mars 2023 \cite{hcpENNVM}. Le classeur local contient les tabulations agrégées. La nomenclature COICOP Maroc 2022 organise les postes par fonction de consommation. La dépense annuelle moyenne par personne (DAMP) est exprimée en dirhams courants par personne et par an. Le coefficient budgétaire (CB) représente la part d'un poste dans le budget publié. Les valeurs constituent une coupe agrégée sur la période de l'enquête.

Sur les 24 feuilles inspectées, le tableau SD~12 de la feuille « Partie3-Structure de la DAMP » fournit les groupes de dépenses par quintile et milieu. Il commence à la ligne 288. Les blocs de données occupent les lignes 294--308, 315--329 et 336--350. Le script \texttt{clean\_households.py} vérifie les trois agrégats avec SD~01, qui publie les mêmes coefficients et arrondit davantage la DAMP. La feuille « Considérations importantes » donne les seuils de quintile pour chaque périmètre. Les feuilles DAMPG01 à DAMPG14 détaillent les catégories utilisées pour examiner la correspondance.

Les 18 profils comprennent un ensemble et cinq quintiles dans chacun des périmètres national, urbain et rural. Les quintiles sont définis selon la \emph{dépense annuelle par personne}. Chaque classement utilise des seuils propres à son périmètre. La comparaison de C1 urbain à C1 rural porte ainsi sur une position relative au sein de deux distributions. Les 18 profils ne sont pas 18 observations individuelles, et leurs ensembles et quintiles se recouvrent en population. Les 14 coefficients de chaque profil totalisent de 99,8 à 100,1\,\%. Sous l'hypothèse d'un arrondi au dixième, une borne maximale de $14\times0,05=0,70$ point s'applique à l'écart de somme. Les coefficients publiés ont été conservés.

\section{Chaîne d'analyse}
Le notebook 01 sert de trame exploratoire. Le notebook 02 calcule les statistiques à partir du CSV IPC. Le notebook 03 documente la couverture des profils et l'audit des correspondances. Le notebook 04 calcule les trajectoires historiques et les chocs hypothétiques. L'annexe \ref{ann:trace} relie les tableaux et figures aux fichiers du projet.

\chapter{Analyse descriptive de l'inflation}
\section{Indicateurs et statistiques}
Si $I_t$ désigne le niveau de l'IPC au mois $t$, la variation mensuelle et le glissement annuel sont respectivement
\[
  \pi_t^{(1)}=100\left(\frac{I_t}{I_{t-1}}-1\right),
  \qquad
  \pi_t^{(12)}=100\left(\frac{I_t}{I_{t-12}}-1\right).
\]
Les premières dates restent vides lorsque le mois de comparaison précède janvier 2017. Le taux d'inflation moyen annuel d'une année civile complète $a$ vaut $100(\overline I_a/\overline I_{a-1}-1)$, où $\overline I_a$ est la moyenne de ses douze niveaux mensuels. Le glissement de décembre compare uniquement deux mois de décembre. La moyenne des glissements mensuels sur douze mois constitue encore un autre résumé.

Les statistiques du tableau \ref{tab:descriptif} sont calculées sur les observations disponibles. La variance corrigée est $s^2=\sum_i(x_i-\bar x)^2/(n-1)$ et l'écart-type vaut $\sqrt{s^2}$. Ces mesures décrivent la période étudiée. Elles ne constituent pas une estimation d'incertitude sur les profils.
\begin{table}[!htbp]
\centering\small
\caption{Statistiques descriptives de l'IPC général et de ses taux}\label{tab:descriptif}
\begin{tabular}{lrrr}
\hline
Statistique & IPC (points) & Mensuel (\%) & Gliss. annuel (\%) \\
\hline
Observations & 118 & 117 & 106 \\
Moyenne & 109{,}49 & 0{,}18 & 2{,}12 \\
Médiane & 105{,}45 & 0{,}19 & 1{,}19 \\
Minimum & 99{,}30 & -1{,}00 & -0{,}75 \\
Maximum & 121{,}80 & 1{,}79 & 10{,}08 \\
Variance & 67{,}21 & 0{,}31 & 6{,}90 \\
Écart-type & 8{,}20 & 0{,}56 & 2{,}63 \\
\hline
\end{tabular}
\sourcetable{Notebook 02, d'après l'IPC du HCP \cite{hcpIndicateurs}, de janvier 2017 à octobre 2026. La variance s'exprime dans le carré de l'unité de la série.}
\end{table}

\begin{table}[!htbp]
\centering\small
\caption{IPC moyen observé et inflation moyenne annuelle}\label{tab:annuel}
\begin{tabular}{rrrr}
\hline
Année & Mois disponibles & IPC moyen & Taux annuel (\%) \\
\hline
2017 & 12 & 99{,}98 & -- \\
2018 & 12 & 101{,}58 & 1{,}60 \\
2019 & 12 & 101{,}83 & 0{,}25 \\
2020 & 12 & 102{,}47 & 0{,}62 \\
2021 & 12 & 103{,}90 & 1{,}40 \\
2022 & 12 & 110{,}82 & 6{,}66 \\
2023 & 12 & 117{,}57 & 6{,}09 \\
2024 & 12 & 118{,}73 & 0{,}99 \\
2025 & 12 & 119{,}56 & 0{,}70 \\
2026 & 10 & 120{,}23 & -- \\
\hline
\end{tabular}
\sourcetable{Notebook 02, d'après l'IPC du HCP \cite{hcpIndicateurs}, base 100 en 2017. Le taux de 2017 exige 2016. Pour 2026, l'IPC moyen couvre janvier à octobre.}
\end{table}


\section{Évolution générale et pics}
L'IPC général passe de 99,40 en janvier 2017 à 121,80 en octobre 2026. La plus forte hausse mensuelle est de 1,79\,\% en mars 2022. Le glissement annuel culmine à 10,08\,\% en février 2023, avec un taux déjà élevé en janvier 2023 (8,91\,\%). Les taux moyens des années civiles 2022 et 2023 atteignent respectivement 6,66\,\% et 6,09\,\%. Sur janvier--octobre 2026, l'IPC général moyen dépasse celui de janvier--octobre 2025 de 0,43\,\%. L'année 2026 est incomplète dans le fichier.
\figureprojet{02_ipc_general.png}{IPC général du Maroc, base 100 en 2017. Source : HCP, séries mensuelles.}{fig:ipc}
\figureprojet{02_inflation_generale.png}{Variations mensuelles et glissements annuels de l'IPC général. Calculs du projet à partir du HCP.}{fig:taux}

\section{Comparaison des postes}
Au pic de février 2023, le glissement de l'agrégat alimentaire atteint 20,08\,\%, contre 3,58\,\% pour les produits non alimentaires. L'écart est de 16,49 points, calculé avant arrondi. Parmi les douze divisions, la moyenne des glissements disponibles de janvier 2018 à octobre 2026 varie de 0,18\,\% pour la santé à 4,68\,\% pour les boissons alcoolisées, tabac et stupéfiants. Cette moyenne compare le rythme des prix entre divisions. Leur contribution à l'IPC général dépend également de leurs poids.
\figureprojet{02_alimentation_non_alimentaire.png}{Agrégats alimentaire et non alimentaire. Les niveaux et les glissements annuels sont représentés sur deux axes distincts.}{fig:alim}
\figureprojet{02_inflation_divisions.png}{Moyenne des glissements annuels disponibles par division, de janvier 2018 à octobre 2026.}{fig:divisions}
\clearpage

\chapter{Correspondances et modèle à pondérations fixes}
\section{Périmètres et couverture}
La correspondance confronte les sous-groupes de l'ENNVM aux groupes et classes de l'IPC détaillés par le HCP \cite{hcpIPC,hcpAgadir}. Le rapport régional d'Agadir étaye la lecture des catégories. Les séries utilisées pour les calculs restent les séries nationales du CSV, collectées en milieu urbain. L'annexe \ref{ann:correspondance} documente les quatorze rapprochements. Le périmètre prudent retient 01, 03, 05 et 10. Le périmètre élargi ajoute les huit correspondances partielles, soit les groupes 01 à 11 et 13. Les groupes 12 (assurances et services financiers) et 14 (paiements fiscaux, transferts et remboursements) restent hors des deux périmètres. Le premier manque d'une série autonome adéquate. Pour le second, les flux recensés ne correspondent pas tous à un prix de consommation.

La couverture des 18 profils apparaît au tableau \ref{tab:profils}. L'ensemble rural consacre 45,0\,\% du budget aux aliments et boissons non alcoolisées, contre 30,1\,\% pour l'ensemble urbain. La part du groupe 01 va de 50,2\,\% pour C1 rural à 24,6\,\% pour C5 urbain. Ces différences décrivent les budgets agrégés et éclairent le mécanisme du choc alimentaire.
\begin{table}[!htbp]
\centering\small
\renewcommand{\arraystretch}{1.10}
\caption{Dépenses annuelles, alimentation et couverture des 18 profils (2022--2023)}\label{tab:profils}
\begin{tabular}{llrrrr}
\hline
Milieu & Classe & DAMP (MAD) & Groupe 01 (\%) & Prudent (\%) & Élargi (\%) \\
\hline
National & C1 & 6\,943{,}10 & 46{,}0 & 51{,}9 & 98{,}0 \\
National & C2 & 10\,848{,}12 & 42{,}6 & 48{,}8 & 97{,}1 \\
National & C3 & 14\,792{,}97 & 39{,}6 & 46{,}2 & 95{,}9 \\
National & C4 & 21\,071{,}63 & 36{,}0 & 43{,}0 & 94{,}5 \\
National & C5 & 49\,633{,}55 & 26{,}7 & 36{,}1 & 90{,}1 \\
National & \textbf{Ensemble} & 20\,658{,}19 & 33{,}4 & 41{,}4 & 93{,}2 \\
\noalign{\vskip 3pt}
Urbain & C1 & 8\,599{,}46 & 39{,}1 & 45{,}0 & 97{,}7 \\
Urbain & C2 & 13\,171{,}00 & 36{,}9 & 43{,}5 & 96{,}2 \\
Urbain & C3 & 17\,835{,}66 & 35{,}2 & 42{,}3 & 95{,}1 \\
Urbain & C4 & 25\,414{,}87 & 32{,}6 & 40{,}4 & 93{,}4 \\
Urbain & C5 & 59\,472{,}69 & 24{,}6 & 34{,}6 & 89{,}5 \\
Urbain & \textbf{Ensemble} & 24\,898{,}22 & 30{,}1 & 38{,}6 & 92{,}4 \\
\noalign{\vskip 3pt}
Rural & C1 & 5\,575{,}70 & 50{,}2 & 55{,}9 & 98{,}3 \\
Rural & C2 & 8\,358{,}93 & 48{,}9 & 55{,}1 & 97{,}8 \\
Rural & C3 & 10\,827{,}81 & 48{,}0 & 54{,}0 & 97{,}0 \\
Rural & C4 & 14\,244{,}14 & 46{,}4 & 52{,}6 & 95{,}8 \\
Rural & C5 & 26\,047{,}17 & 40{,}6 & 46{,}5 & 92{,}7 \\
Rural & \textbf{Ensemble} & 13\,009{,}92 & 45{,}0 & 51{,}0 & 95{,}2 \\
\hline
\end{tabular}
\sourcetable{HCP, ENNVM 2022--2023, tableau SD~12 \cite{hcpENNVM}. DAMP : dépense annuelle moyenne par personne. C1 à C5 : quintiles de dépenses du milieu. Prudent et élargi : parts du budget couvertes.}
\end{table}


\section{Coût historique du sous-panier}
Pour un profil $p$, un groupe $k$ et un ensemble de postes $S$, le coefficient publié $CB_{p,k}$ est exprimé en pourcentage du budget total. Les poids renormalisés et l'indice de coût du sous-panier sont définis par
\[
 C_p^S=\sum_{k\in S}CB_{p,k},\qquad w_{p,k}^S=\frac{CB_{p,k}}{C_p^S},\qquad
 P_{p,t}^S=100\sum_{k\in S}w_{p,k}^S\frac{I_{k,t}}{I_{k,b}} .
\]
Les poids $w_{p,k}^S$ totalisent 1. Par exemple, le groupe 01 représente 33,4\,\% du budget de l'ensemble national, soit 80,7\,\% du sous-panier prudent et 35,8\,\% du sous-panier élargi. Le mois de référence $b$ est septembre 2022, proche du milieu de la collecte. Ce repère fixe la valeur initiale à 100. Les coefficients décrivent l'ensemble de l'enquête et leur emploi à toute autre date implique un panier de composition constante.

La variation depuis la référence est $P_{p,t}^S-100$, en pourcentage du coût initial du sous-panier. Sa contribution au budget total initial, en points, est
\[
 A_{p,t}^S=\frac{C_p^S}{100}(P_{p,t}^S-100)
     =\sum_{k\in S}CB_{p,k}\left(\frac{I_{k,t}}{I_{k,b}}-1\right).
\]
Les variations à un et douze mois comparent $P_{p,t}^S$ à $P_{p,t-h}^S$. La contribution n'est calculée ici que depuis la référence. Pour un autre horizon, la part du sous-panier au début de l'intervalle pourrait différer.

\section{Chocs hypothétiques}
Avec un choc de prix $\delta_k$ en fraction et un budget initial hypothétique $B=5\,000$ MAD, la contribution, le taux du sous-panier et le supplément s'écrivent
\[
 A_p^S=\sum_{k\in S}CB_{p,k}\delta_k,\qquad
 \Delta_p^S=100\frac{A_p^S}{C_p^S},\qquad
 M_p^S=B\frac{A_p^S}{100}.
\]
Le choc uniforme fixe chaque $\delta_k$ à 0,05. Dans le choc combiné, le groupe 01 augmente ensuite de 10\,\% sur son prix déjà majoré, soit $(1{,}05)(1{,}10)-1=0{,}155$. Les autres postes restent à +5\,\%. La composante alimentaire supplémentaire vaut donc 10,5 points du prix alimentaire initial. Seuls les postes couverts contribuent au supplément. L'illustration suppose le même budget de départ pour tous les profils.

\chapter{Résultats des simulations}
\section{Trajectoires aux prix historiques}
La figure \ref{fig:historique} présente les 18 profils dans les deux périmètres. Les trajectoires utilisent les indices de prix observés du HCP et les coefficients fixes de l'ENNVM. Les variations diffèrent selon le périmètre et la composition des budgets. Entre septembre 2022 et octobre 2026, les hausses de coût sont comprises entre 8,97 et 9,45\,\% pour le périmètre prudent, et entre 7,60 et 8,34\,\% pour l'élargi. Pour l'ensemble national, les hausses sont de 9,22 et 7,85\,\%. Les contributions correspondantes au budget initial sont respectivement de 3,82 et 7,31 points. Le sous-panier prudent augmente davantage en proportion de son propre coût initial. Le périmètre élargi pèse davantage dans le budget total, ce qui conduit à une contribution supérieure.

Les variations d'octobre 2026 sur un mois ou sur douze mois sont disponibles dans le CSV historique. Le tableau \ref{tab:historique} résume les 18 profils à cette date par rapport à la référence. Les trajectoires antérieures à septembre 2022 appliquent rétrospectivement les structures budgétaires de l'enquête.
\figureprojet[1.10\linewidth]{04_cout_panier_historique.png}{Indices de coût par milieu et périmètre, de janvier 2017 à octobre 2026 (septembre 2022 = 100). C1 à C5 désignent les quintiles de dépenses de chaque milieu. Certaines courbes se superposent. Source : notebook 04, d'après le HCP \cite{hcpIndicateurs,hcpENNVM}.}{fig:historique}
\clearpage
\begin{table}[!htbp]
\centering\small
\renewcommand{\arraystretch}{1.10}
\caption{Variation du coût et contribution au budget initial en octobre 2026}\label{tab:historique}
\begin{tabular}{llrrrr}
\hline
Milieu & Classe & Prudent (\%) & Prudent (pts) & Élargi (\%) & Élargi (pts) \\
\hline
National & C1 & 9{,}42 & 4{,}89 & 8{,}21 & 8{,}05 \\
National & C2 & 9{,}38 & 4{,}58 & 8{,}16 & 7{,}92 \\
National & C3 & 9{,}36 & 4{,}32 & 8{,}05 & 7{,}72 \\
National & C4 & 9{,}31 & 4{,}00 & 7{,}90 & 7{,}46 \\
National & C5 & 9{,}04 & 3{,}26 & 7{,}65 & 6{,}89 \\
National & \textbf{Ensemble} & 9{,}22 & 3{,}82 & 7{,}85 & 7{,}31 \\
\noalign{\vskip 3pt}
Urbain & C1 & 9{,}39 & 4{,}23 & 8{,}07 & 7{,}89 \\
Urbain & C2 & 9{,}35 & 4{,}07 & 7{,}99 & 7{,}69 \\
Urbain & C3 & 9{,}31 & 3{,}94 & 7{,}93 & 7{,}54 \\
Urbain & C4 & 9{,}26 & 3{,}74 & 7{,}81 & 7{,}29 \\
Urbain & C5 & 8{,}97 & 3{,}10 & 7{,}60 & 6{,}80 \\
Urbain & \textbf{Ensemble} & 9{,}17 & 3{,}54 & 7{,}77 & 7{,}18 \\
\noalign{\vskip 3pt}
Rural & C1 & 9{,}45 & 5{,}28 & 8{,}34 & 8{,}20 \\
Rural & C2 & 9{,}41 & 5{,}18 & 8{,}26 & 8{,}08 \\
Rural & C3 & 9{,}41 & 5{,}08 & 8{,}29 & 8{,}04 \\
Rural & C4 & 9{,}39 & 4{,}94 & 8{,}22 & 7{,}88 \\
Rural & C5 & 9{,}34 & 4{,}34 & 8{,}00 & 7{,}41 \\
Rural & \textbf{Ensemble} & 9{,}38 & 4{,}78 & 8{,}16 & 7{,}77 \\
\hline
\end{tabular}
\sourcetable{Notebook 04, d'après le HCP \cite{hcpIndicateurs,hcpENNVM}. Référence : septembre 2022. Les points rapportent la variation des postes couverts au budget initial complet.}
\end{table}


\section{Choc uniforme de \texorpdfstring{5\,\%}{5 \%}}
Le coût du sous-panier augmente de 5\,\% dans les 36 couples profil--périmètre. Dans l'illustration de 5\,000 MAD, le supplément sur les postes couverts varie entre 86,50 et 139,75 MAD dans le périmètre prudent, et entre 223,75 et 245,75 MAD dans l'élargi. À hausse uniforme donnée, cet écart reflète la proportion du budget couverte. Une hausse de 5\,\% sur le budget entier de 5\,000 MAD donnerait 250 MAD.

\section{Choc uniforme suivi du choc alimentaire}
Le supplément du scénario combiné varie entre 215,65 et 403,30 MAD dans le périmètre prudent, et entre 352,90 et 509,30 MAD dans l'élargi. Pour l'ensemble national, le coût du sous-panier prudent augmente de 13,47\,\% (+278,85 MAD et +5,58 points du budget initial), contre 8,76\,\% dans l'élargi (+408,35 MAD et +8,17 points). L'ajout du choc alimentaire représente 175,35 MAD dans les deux périmètres de ce profil, car l'alimentation reçoit le même coefficient d'origine. Parmi les profils, les plus fortes charges alimentaires supplémentaires correspondent aux parts alimentaires initiales les plus élevées. Ainsi, C1 rural consacre 50,2\,\% de son budget au groupe 01, ce qui correspond à 263,55 MAD de supplément alimentaire marginal sur 5\,000 MAD.

Les tableaux \ref{tab:chocs-prudent} et \ref{tab:chocs-elargi} présentent les montants et taux par profil. Les montants illustrent les dépenses supplémentaires nécessaires au maintien des quantités couvertes. Les dépenses réellement engagées dépendraient notamment des substitutions possibles, des prix effectivement rencontrés et des ressources du ménage.
Les contributions au budget initial figurent en annexe~\ref{ann:chocs}. \figureprojet[1.10\linewidth]{04_chocs_hypothetiques.png}{Suppléments en MAD pour un budget initial hypothétique de 5\,000 MAD. Les barres distinguent la hausse uniforme et le supplément alimentaire. « Total » désigne l'ensemble du milieu. Source : notebook 04, pondérations ENNVM \cite{hcpENNVM}.}{fig:chocs}
\clearpage
\begin{table}[!htbp]
\centering\small
\caption{Suppléments sur 5\,000 MAD, périmètre prudent}\label{tab:chocs-prudent}
\begin{tabular}{llrrrr}
\hline
Milieu & Classe & Couvert (\%) & Uniforme (MAD) & Combiné (\%) & Combiné (MAD) \\
\hline
National & C1 & 51{,}9 & 129{,}75 & 14{,}31 & 371{,}25 \\
National & C2 & 48{,}8 & 122{,}00 & 14{,}17 & 345{,}65 \\
National & C3 & 46{,}2 & 115{,}50 & 14{,}00 & 323{,}40 \\
National & C4 & 43{,}0 & 107{,}50 & 13{,}79 & 296{,}50 \\
National & C5 & 36{,}1 & 90{,}25 & 12{,}77 & 230{,}43 \\
National & Ensemble & 41{,}4 & 103{,}50 & 13{,}47 & 278{,}85 \\
Urbain & C1 & 45{,}0 & 112{,}50 & 14{,}12 & 317{,}77 \\
Urbain & C2 & 43{,}5 & 108{,}75 & 13{,}91 & 302{,}48 \\
Urbain & C3 & 42{,}3 & 105{,}75 & 13{,}74 & 290{,}55 \\
Urbain & C4 & 40{,}4 & 101{,}00 & 13{,}47 & 272{,}15 \\
Urbain & C5 & 34{,}6 & 86{,}50 & 12{,}47 & 215{,}65 \\
Urbain & Ensemble & 38{,}6 & 96{,}50 & 13{,}19 & 254{,}53 \\
Rural & C1 & 55{,}9 & 139{,}75 & 14{,}43 & 403{,}30 \\
Rural & C2 & 55{,}1 & 137{,}75 & 14{,}32 & 394{,}48 \\
Rural & C3 & 54{,}0 & 135{,}00 & 14{,}33 & 387{,}00 \\
Rural & C4 & 52{,}6 & 131{,}50 & 14{,}26 & 375{,}10 \\
Rural & C5 & 46{,}5 & 116{,}25 & 14{,}17 & 329{,}40 \\
Rural & Ensemble & 51{,}0 & 127{,}50 & 14{,}26 & 363{,}75 \\
\hline
\end{tabular}
\sourcetable{Simulations hypothétiques. Le taux uniforme vaut 5\,\% dans tous les cas.}
\end{table}

\begin{table}[!htbp]
\centering\small
\caption{Suppléments sur 5\,000 MAD, périmètre élargi}\label{tab:chocs-elargi}
\begin{tabular}{llrrrr}
\hline
Milieu & Classe & Couvert (\%) & Uniforme (MAD) & Combiné (\%) & Combiné (MAD) \\
\hline
National & C1 & 98{,}0 & 245{,}00 & 9{,}93 & 486{,}50 \\
National & C2 & 97{,}1 & 242{,}75 & 9{,}61 & 466{,}40 \\
National & C3 & 95{,}9 & 239{,}75 & 9{,}34 & 447{,}65 \\
National & C4 & 94{,}5 & 236{,}25 & 9{,}00 & 425{,}25 \\
National & C5 & 90{,}1 & 225{,}25 & 8{,}11 & 365{,}43 \\
National & Ensemble & 93{,}2 & 233{,}00 & 8{,}76 & 408{,}35 \\
Urbain & C1 & 97{,}7 & 244{,}25 & 9{,}20 & 449{,}52 \\
Urbain & C2 & 96{,}2 & 240{,}50 & 9{,}03 & 434{,}23 \\
Urbain & C3 & 95{,}1 & 237{,}75 & 8{,}89 & 422{,}55 \\
Urbain & C4 & 93{,}4 & 233{,}50 & 8{,}66 & 404{,}65 \\
Urbain & C5 & 89{,}5 & 223{,}75 & 7{,}89 & 352{,}90 \\
Urbain & Ensemble & 92{,}4 & 231{,}00 & 8{,}42 & 389{,}02 \\
Rural & C1 & 98{,}3 & 245{,}75 & 10{,}36 & 509{,}30 \\
Rural & C2 & 97{,}8 & 244{,}50 & 10{,}25 & 501{,}23 \\
Rural & C3 & 97{,}0 & 242{,}50 & 10{,}20 & 494{,}50 \\
Rural & C4 & 95{,}8 & 239{,}50 & 10{,}09 & 483{,}10 \\
Rural & C5 & 92{,}7 & 231{,}75 & 9{,}60 & 444{,}90 \\
Rural & Ensemble & 95{,}2 & 238{,}00 & 9{,}96 & 474{,}25 \\
\hline
\end{tabular}
\sourcetable{Simulations hypothétiques. Le taux uniforme vaut 5\,\% dans tous les cas.}
\end{table}


\clearpage
\section{Pouvoir d'achat du sous-panier et ressources nominales hypothétiques}
\label{sec:pouvoir-achat}
\subsection{Hypothèses, mesure et seuil de maintien}

Les quatre hypothèses de croissance des ressources nominales sont 0\,\%, +3\,\%, +5\,\% et +10\,\%. Elles s'appliquent au montant effectivement affecté au sous-panier dans l'expérience de simulation. Elles ne proviennent d'aucune série de revenus observés par le HCP. Pour les comparaisons aux prix historiques, chaque hypothèse désigne une variation \emph{totale} entre septembre 2022 et la date considérée. Ces taux ne sont ni annualisés ni composés mois par mois. Pour les deux chocs hypothétiques, ils s'appliquent une seule fois entre l'avant et l'après-choc.

Soient $R_0>0$ les ressources nominales initialement affectées au sous-panier et $C_0>0$ son coût initial. Le nombre d'unités accessibles passe de $R_0/C_0$ à $R_0(1+g_R)/[C_0(1+\pi_h)]$. Pour des taux exprimés en fractions, avec $1+\pi_h>0$, sa variation relative exacte vaut
\begin{equation}
 \Delta_{\mathrm{PA}}
 = \frac{1+g_R}{1+\pi_h}-1
 = \frac{g_R-\pi_h}{1+\pi_h}.
 \label{eq:pouvoir-achat}
\end{equation}
Le résultat est multiplié par 100 dans les tableaux. Le seuil de maintien exact est $g_R^\star=\pi_h$. En dessous, les ressources achètent moins d'unités du sous-panier, et au-dessus elles en achètent davantage. Si le coût diminue, le seuil est négatif. La couverture est reportée pour situer le sous-panier dans le budget total, sans intervenir comme multiplicateur du taux de pouvoir d'achat.

\subsection{Comparaison d'octobre 2026 à septembre 2022}

Pour l'ensemble national, une hausse nominale de 5\,\% donne une variation de $-3{,}86$\,\% dans le périmètre prudent et de $-2{,}64$\,\% dans l'élargi. Le maintien exigerait des hausses nominales de 9,22\,\% et 7,85\,\%, respectivement. Au même taux de ressources, les ensembles urbain et rural affichent respectivement $-3{,}82$\,\% et $-4{,}01$\,\% dans le prudent, puis $-2{,}57$\,\% et $-2{,}93$\,\% dans l'élargi.

Dans chaque milieu, C1 subit une contraction plus forte que C5 pour cette comparaison. Dans le périmètre prudent, le contraste atteint par exemple $-4{,}02$\,\% contre $-3{,}64$\,\% en milieu urbain avec +5\,\% de ressources. Les quintiles sont définis sur les dépenses par personne avec des seuils propres à chaque milieu. Les tableaux \ref{tab:pa-h-prudent} et \ref{tab:pa-h-elargi} détaillent les quatre hypothèses. En octobre 2026, +5\,\% reste inférieur au seuil pour les 18 profils dans les deux périmètres. +10\,\% dépasse tous les seuils de cette comparaison historique.
\figureprojet{05_pouvoir_achat_historique.png}{Variation du nombre d'unités accessibles des sous-paniers en octobre 2026, en référence à septembre 2022. Coûts calculés à partir des prix historiques et croissance nominale hypothétique. Source : notebook 05.}{fig:pa-historique}
\clearpage
\begin{table}[!htbp]
\centering\small
\caption{Octobre 2026 par rapport à septembre 2022, périmètre prudent}\label{tab:pa-h-prudent}
\begin{tabular}{llrrrrrr}
\hline
Milieu & Classe & Couvert & Seuil & 0\,\% & +3\,\% & +5\,\% & +10\,\% \\
\hline
National & C1 & 51{,}9 & 9{,}42 & -8{,}61 & -5{,}87 & -4{,}04 & 0{,}53 \\
National & C2 & 48{,}8 & 9{,}38 & -8{,}58 & -5{,}83 & -4{,}01 & 0{,}57 \\
National & C3 & 46{,}2 & 9{,}36 & -8{,}56 & -5{,}81 & -3{,}99 & 0{,}59 \\
National & C4 & 43{,}0 & 9{,}31 & -8{,}52 & -5{,}77 & -3{,}94 & 0{,}63 \\
National & C5 & 36{,}1 & 9{,}04 & -8{,}29 & -5{,}54 & -3{,}71 & 0{,}88 \\
National & Ensemble & 41{,}4 & 9{,}22 & -8{,}44 & -5{,}70 & -3{,}86 & 0{,}71 \\
\noalign{\vskip 2pt}
Urbain & C1 & 45{,}0 & 9{,}39 & -8{,}59 & -5{,}84 & -4{,}01 & 0{,}56 \\
Urbain & C2 & 43{,}5 & 9{,}35 & -8{,}55 & -5{,}80 & -3{,}97 & 0{,}60 \\
Urbain & C3 & 42{,}3 & 9{,}31 & -8{,}52 & -5{,}77 & -3{,}94 & 0{,}63 \\
Urbain & C4 & 40{,}4 & 9{,}26 & -8{,}48 & -5{,}73 & -3{,}90 & 0{,}67 \\
Urbain & C5 & 34{,}6 & 8{,}97 & -8{,}23 & -5{,}48 & -3{,}64 & 0{,}95 \\
Urbain & Ensemble & 38{,}6 & 9{,}17 & -8{,}40 & -5{,}65 & -3{,}82 & 0{,}76 \\
\noalign{\vskip 2pt}
Rural & C1 & 55{,}9 & 9{,}45 & -8{,}63 & -5{,}89 & -4{,}06 & 0{,}50 \\
Rural & C2 & 55{,}1 & 9{,}41 & -8{,}60 & -5{,}86 & -4{,}03 & 0{,}54 \\
Rural & C3 & 54{,}0 & 9{,}41 & -8{,}60 & -5{,}86 & -4{,}03 & 0{,}54 \\
Rural & C4 & 52{,}6 & 9{,}39 & -8{,}58 & -5{,}84 & -4{,}01 & 0{,}56 \\
Rural & C5 & 46{,}5 & 9{,}34 & -8{,}54 & -5{,}80 & -3{,}97 & 0{,}60 \\
Rural & Ensemble & 51{,}0 & 9{,}38 & -8{,}58 & -5{,}83 & -4{,}01 & 0{,}57 \\
\hline
\end{tabular}
\sourcetable{Notebook 05. Couvert : part du budget initial en pourcentage. Seuil : croissance nominale de maintien en pourcentage. Les quatre dernières colonnes donnent la variation des unités accessibles (en \%).}
\end{table}

\begin{table}[!htbp]
\centering\small
\renewcommand{\arraystretch}{1.10}
\caption{Octobre 2026 par rapport à septembre 2022, périmètre élargi}\label{tab:pa-h-elargi}
\begin{tabular}{llrrrrrr}
\hline
Milieu & Classe & Couverture & Seuil & 0\,\% & +3\,\% & +5\,\% & +10\,\% \\
\hline
National & C1 & 98{,}0 & 8{,}21 & -7{,}59 & -4{,}82 & -2{,}97 & 1{,}65 \\
National & C2 & 97{,}1 & 8{,}16 & -7{,}54 & -4{,}77 & -2{,}92 & 1{,}70 \\
National & C3 & 95{,}9 & 8{,}05 & -7{,}45 & -4{,}68 & -2{,}82 & 1{,}80 \\
National & C4 & 94{,}5 & 7{,}90 & -7{,}32 & -4{,}54 & -2{,}68 & 1{,}95 \\
National & C5 & 90{,}1 & 7{,}65 & -7{,}11 & -4{,}32 & -2{,}46 & 2{,}18 \\
National & \textbf{Ensemble} & 93{,}2 & 7{,}85 & -7{,}27 & -4{,}49 & -2{,}64 & 2{,}00 \\
\noalign{\vskip 3pt}
Urbain & C1 & 97{,}7 & 8{,}07 & -7{,}47 & -4{,}69 & -2{,}84 & 1{,}78 \\
Urbain & C2 & 96{,}2 & 7{,}99 & -7{,}40 & -4{,}62 & -2{,}77 & 1{,}86 \\
Urbain & C3 & 95{,}1 & 7{,}93 & -7{,}34 & -4{,}56 & -2{,}71 & 1{,}92 \\
Urbain & C4 & 93{,}4 & 7{,}81 & -7{,}24 & -4{,}46 & -2{,}60 & 2{,}03 \\
Urbain & C5 & 89{,}5 & 7{,}60 & -7{,}06 & -4{,}27 & -2{,}41 & 2{,}23 \\
Urbain & \textbf{Ensemble} & 92{,}4 & 7{,}77 & -7{,}21 & -4{,}43 & -2{,}57 & 2{,}07 \\
\noalign{\vskip 3pt}
Rural & C1 & 98{,}3 & 8{,}34 & -7{,}70 & -4{,}93 & -3{,}09 & 1{,}53 \\
Rural & C2 & 97{,}8 & 8{,}26 & -7{,}63 & -4{,}86 & -3{,}01 & 1{,}61 \\
Rural & C3 & 97{,}0 & 8{,}29 & -7{,}66 & -4{,}89 & -3{,}04 & 1{,}58 \\
Rural & C4 & 95{,}8 & 8{,}22 & -7{,}60 & -4{,}83 & -2{,}98 & 1{,}64 \\
Rural & C5 & 92{,}7 & 8{,}00 & -7{,}40 & -4{,}63 & -2{,}78 & 1{,}85 \\
Rural & \textbf{Ensemble} & 95{,}2 & 8{,}16 & -7{,}55 & -4{,}77 & -2{,}92 & 1{,}70 \\
\hline
\end{tabular}
\sourcetable{Notebook 05, d'après le HCP \cite{hcpIndicateurs,hcpENNVM}. Couverture et seuil de maintien en \%. Les quatre dernières colonnes donnent la variation des unités accessibles (en \%) par hypothèse nominale.}
\end{table}

\clearpage

\subsection{Chocs hypothétiques et comparaison des périmètres}

Pour la hausse uniforme des prix de 5\,\%, le seuil vaut 5\,\% pour chacun des 36 couples profil--périmètre. Les scénarios nominaux 0\,\%, +3\,\%, +5\,\% et +10\,\% donnent respectivement $-4{,}76$\,\%, $-1{,}90$\,\%, 0\,\% et $+4{,}76$\,\% en unités du sous-panier.

Quand la hausse uniforme est suivie du choc alimentaire supplémentaire, le seuil varie de 12,47\,\% à 14,43\,\% dans le prudent et de 7,89\,\% à 10,36\,\% dans l'élargi. Une hausse nominale de 10\,\% suffit pour 14 des 18 profils dans l'élargi. Les quatre exceptions sont les quintiles ruraux C1 à C4. Elle reste inférieure au seuil pour tous les profils du prudent. Le poste alimentaire représente une part plus élevée du sous-panier prudent, ce qui amplifie l'effet relatif du choc. Les taux des deux périmètres portent sur des compositions et des couvertures différentes.

La figure \ref{fig:pa-chocs} compare les deux chocs. L'annexe \ref{ann:pa-chocs} détaille les quatre variations et le seuil pour les 18 profils après le choc combiné.
\figureprojet{05_pouvoir_achat_chocs.png}{Variations du pouvoir d'achat des sous-paniers sous les chocs de prix et les hypothèses de croissance des ressources. Source : notebook 05.}{fig:pa-chocs}
\clearpage

\subsection{Limites d'interprétation}

Les résultats mesurent une capacité conditionnelle à financer un sous-panier. Les profils sont des agrégats de dépenses et leurs quintiles ne renseignent pas sur le classement des revenus. Le modèle ne dispose ni de trajectoires de revenus effectivement observées ni de quantités suivies permettant d'étudier les substitutions. Il conserve les pondérations de l'ENNVM 2022--2023 à toutes les dates et suppose une évolution proportionnelle des ressources affectées aux postes couverts. Les profils ruraux reçoivent les mêmes variations de prix à collecte urbaine que les profils urbains \cite{hcpIPC}. Les écarts simulés par milieu s'interprètent donc sous cette hypothèse de prix communs. Les résultats ne s'étendent pas aux postes exclus ni au pouvoir d'achat de l'ensemble du budget.


\chapter{Discussion critique}
\section{Arbitrage entre couverture et précision}
Le scénario prudent réduit le risque d'appariement en s'appuyant sur quatre groupes dont les sous-postes concordent mieux. La couverture se situe entre 34,6 et 55,9\,\%. Le groupe alimentaire domine ce sous-panier dans plusieurs profils, si bien que la variation estimée peut être particulièrement sensible à son évolution. Le scénario élargi couvre entre 89,5 et 98,3\,\%, à l'aide de correspondances partielles. Le logement comprend des loyers imputés dans l'ENNVM alors que le suivi IPC des loyers repose sur des locataires. Les frontières du transport, des communications et des loisirs diffèrent aussi entre nomenclatures. Les écarts entre les deux scénarios combinent donc l'effet des postes ajoutés et la redistribution des poids. Ils constituent un test de sensibilité au périmètre.

\section{Collecte urbaine et comparabilité des profils}
En l'absence de séries rurales comparables dans les fichiers analysés, le même vecteur de prix est appliqué à tous les profils. L'écart simulé entre un profil rural et un profil urbain isole le rôle des structures budgétaires conditionnellement à cette hypothèse. Les prix locaux, l'accès aux marchés, la disponibilité des produits et l'autoconsommation peuvent modifier l'exposition effective. Les seuils de quintiles propres à chaque milieu compliquent en outre les comparaisons à niveau de dépense absolu comparable. La couverture des groupes de dépenses ne mesure pas la représentativité géographique.

\section{Portée des résultats et vérifications complémentaires}
Le modèle tient les poids du panier constants et utilise des indices déjà agrégés par division. Il ne permet pas d'estimer une élasticité de demande ou de quantifier les substitutions. Les séries et profils ne fournissent pas l'information nécessaire pour identifier une relation causale entre inflation et comportement ou chiffrer une perte effective de pouvoir d'achat. Les profils publiés ne permettent pas non plus de calculer des intervalles d'incertitude issus du plan de sondage.

Trois validations ont la priorité pour approfondir l'étude. La première consiste à vérifier les appariements de sous-postes, notamment les loyers, la santé, les reclassements entre divisions et les services divers. La deuxième serait de comparer les résultats à des prix effectivement relevés en milieu rural. La troisième consiste à tester le mois de référence, la stabilité des poids, et l'usage éventuel d'une moyenne des prix sur la collecte. L'enquête ne fournit pas dans les tables exploitées de poids mensuels permettant une actualisation observée.

\clearpage
\section{Analyse de robustesse et sensibilité des résultats}
\label{sec:robustesse}
Le notebook 06 compare les 18 profils agrégés de l'ENNVM dans les périmètres prudent et élargi. Il fait varier le mois de référence et la composition du sous-panier. La couverture initiale est conservée pour chaque profil. Elle varie de 34,6 à 55,9\,\% dans le prudent et de 89,5 à 98,3\,\% dans l'élargi.

\subsection{Redistribution proportionnelle des pondérations}

Soit $w_k$ le poids du poste $k$ en pourcentage du sous-panier, après renormalisation des coefficients publiés. La variante initiale conserve ces poids. Les deux variantes hypothétiques modifient le poids du groupe 01, aliments et boissons non alcoolisées, de $\delta=-5$ ou $\delta=+5$ points de pourcentage. Pour répartir le complément, on conserve les proportions entre tous les autres postes :
\[
 w'_1=w_1+\delta,\qquad
 w'_k=w_k\,\frac{100-w'_1}{100-w_1}\quad(k\ne1).
\]
Comme $\sum_{k\ne 1}w_k=100-w_1$, on obtient
\[
 \sum_k w'_k=w'_1+\frac{100-w'_1}{100-w_1}(100-w_1)=100 .
\]
Les poids sont strictement positifs dans les variantes testées. Cette opération modifie la composition du sous-panier et conserve sa couverture et ses catégories. Elle représente une hypothèse de sensibilité, sans estimation de substitutions ni nouvelle observation de dépenses.

\subsection{Sensibilité historique et horizons de comparaison}

Pour un mois de référence $b$, le coût relatif du sous-panier au mois $t$ vaut
\[
 100\sum_k \frac{w'_k}{100}\frac{I_{k,t}}{I_{k,b}}.
\]
Chaque référence applique de nouveau les poids à son propre point de départ. Il s'agit d'une nouvelle comparaison à poids fixes pour chaque fenêtre. La variation cumulée, en pourcentage, est l'indice obtenu diminué de 100.

Le tableau \ref{tab:robustesse-histoire} et la figure \ref{fig:robustesse-histoire} comparent deux familles de fenêtres. Les trois fenêtres de douze mois ont la même durée et traversent des périodes calendaires différentes. Elles testent la sensibilité au choix de la période. Les fenêtres se terminant toutes en octobre 2026 durent 49, 45 ou 34 mois. Leurs variations cumulées décrivent des durées différentes et ne peuvent donc pas être lues comme une comparaison directe de rythmes. Pour l'ensemble national, avec les poids initiaux, le prudent augmente davantage que l'élargi dans les deux premières fenêtres de douze mois. L'ordre s'inverse de décembre 2023 à décembre 2024, avec +0,73\,\% dans le prudent et +1,34\,\% dans l'élargi. Sur cette dernière fenêtre, donner cinq points supplémentaires à l'alimentation abaisse les taux à +0,70\,\% et +1,28\,\%, respectivement.
\begin{table}[!htbp]
\centering\small
\setlength{\tabcolsep}{4pt}
\caption{Sensibilité de la variation du coût pour l'ensemble national (en \%)}\label{tab:robustesse-histoire}
\begin{tabular}{lllrrrrrr}
\hline
\multicolumn{3}{c}{Fenêtre} & \multicolumn{3}{c}{Prudent} & \multicolumn{3}{c}{Élargi} \\
Début & Fin & Mois & Initial & $-5$ pts & $+5$ pts & Initial & $-5$ pts & $+5$ pts \\
\hline
2022-09 & 2023-09 & 12 & 8{,}72 & 8{,}36 & 9{,}09 & 4{,}60 & 4{,}17 & 5{,}03 \\
2023-01 & 2024-01 & 12 & 3{,}86 & 3{,}76 & 3{,}96 & 2{,}21 & 2{,}05 & 2{,}37 \\
2023-12 & 2024-12 & 12 & 0{,}73 & 0{,}76 & 0{,}70 & 1{,}34 & 1{,}39 & 1{,}28 \\
\hline
2022-09 & 2026-10 & 49 & 9{,}22 & 9{,}07 & 9{,}37 & 7{,}85 & 7{,}69 & 8{,}00 \\
2023-01 & 2026-10 & 45 & 6{,}22 & 6{,}18 & 6{,}26 & 6{,}32 & 6{,}32 & 6{,}33 \\
2023-12 & 2026-10 & 34 & 1{,}40 & 1{,}52 & 1{,}28 & 3{,}56 & 3{,}77 & 3{,}36 \\
\hline
\end{tabular}
\sourcetable{Notebook 06. Les trois premières lignes ont un horizon commun de douze mois. Les trois dernières ont une arrivée commune en octobre 2026.}
\end{table}

\figureprojet{sensibilite_historique.png}{Sensibilité de la variation du coût (en \%) pour l'ensemble national, selon les dates de référence et d'arrivée indiquées. Les points correspondent aux poids initiaux et les segments à l'étendue des variantes alimentaires. Les fenêtres de durée différente figurent dans un panneau distinct. Source : notebook 06, d'après le HCP \cite{hcpIndicateurs,hcpENNVM}.}{fig:robustesse-histoire}
\clearpage

\subsection{Rural C1 et Urbain C5 face au choc alimentaire}

Le choc uniforme augmente de 5\,\% le prix de chaque poste couvert. Le choc combiné augmente ensuite de 10\,\% le prix alimentaire déjà majoré, soit une hausse totale de 15,5\,\% pour ce poste. Le supplément alimentaire marginal est la différence entre les montants du choc combiné et du choc uniforme. Sur le budget illustratif de 5\,000 MAD, il vaut 10,5\,\% du coût initial du poste alimentaire. Le tableau \ref{tab:robustesse-choc} distingue ce montant de la variation du coût de tout le sous-panier.

Avec les coefficients publiés, la part alimentaire du budget total est de 50,2\,\% pour Rural C1 et de 24,6\,\% pour Urbain C5. Le supplément marginal correspondant est de 263,55 MAD contre 129,15 MAD dans les deux périmètres, soit un écart de 134,40 MAD pour le même budget illustratif. Il est identique entre les périmètres avec les poids initiaux parce que le coefficient du groupe 01 dans le budget total est alors le même. Sous les variantes, le poids modifié est défini à l'intérieur du sous-panier et sa conversion en part du budget total dépend de la couverture. La figure \ref{fig:robustesse-choc} illustre cette sensibilité.
\begin{table}[!htbp]
\centering\small
\setlength{\tabcolsep}{4pt}
\caption{Choc combiné et supplément alimentaire marginal pour les deux profils}\label{tab:robustesse-choc}
\begin{tabular}{lllrrrr}
\hline
Périmètre & Profil & Variante & Couv. (\%) & Poids alim. (\%) & Coût (\%) & Marge (MAD) \\
\hline
Prudent & Rural C1 & Initial & 55{,}9 & 89{,}80 & 14{,}43 & 263{,}55 \\
Prudent & Rural C1 & $-5$ pts & 55{,}9 & 84{,}80 & 13{,}90 & 248{,}88 \\
Prudent & Rural C1 & $+5$ pts & 55{,}9 & 94{,}80 & 14{,}95 & 278{,}22 \\
Prudent & Urbain C5 & Initial & 34{,}6 & 71{,}10 & 12{,}47 & 129{,}15 \\
Prudent & Urbain C5 & $-5$ pts & 34{,}6 & 66{,}10 & 11{,}94 & 120{,}07 \\
Prudent & Urbain C5 & $+5$ pts & 34{,}6 & 76{,}10 & 12{,}99 & 138{,}23 \\
\hline
Élargi & Rural C1 & Initial & 98{,}3 & 51{,}07 & 10{,}36 & 263{,}55 \\
Élargi & Rural C1 & $-5$ pts & 98{,}3 & 46{,}07 & 9{,}84 & 237{,}75 \\
Élargi & Rural C1 & $+5$ pts & 98{,}3 & 56{,}07 & 10{,}89 & 289{,}35 \\
Élargi & Urbain C5 & Initial & 89{,}5 & 27{,}49 & 7{,}89 & 129{,}15 \\
Élargi & Urbain C5 & $-5$ pts & 89{,}5 & 22{,}49 & 7{,}36 & 105{,}66 \\
Élargi & Urbain C5 & $+5$ pts & 89{,}5 & 32{,}49 & 8{,}41 & 152{,}64 \\
\hline
\end{tabular}
\sourcetable{Notebook 06. Le poids alimentaire est celui du sous-panier. Le coût est sa variation sous le choc combiné. La marge correspond au seul supplément alimentaire sur un budget initial hypothétique de 5\,000 MAD.}
\end{table}

\figureprojet{choc_alimentaire_rural_urbain.png}{Supplément alimentaire marginal pour Rural C1 et Urbain C5, sur un budget hypothétique de 5\,000 MAD. Source : notebook 06, pondérations ENNVM \cite{hcpENNVM}.}{fig:robustesse-choc}
\clearpage

\subsection{Résultats robustes, résultats sensibles et limites}

\textbf{Résultats robustes dans les variantes testées.} La hausse uniforme de 5\,\% augmente de 5\,\% le coût de chaque sous-panier. Le supplément alimentaire marginal est plus élevé pour Rural C1 que pour Urbain C5 dans les deux périmètres, y compris lorsque la variante la plus basse pour Rural C1 est comparée à la plus haute pour Urbain C5 au sein du même périmètre.

\textbf{Résultats sensibles.} Le classement historique des périmètres varie selon la fenêtre. Dans l'élargi, la hausse du coût de Rural C1 sous le choc combiné va de 9,84 à 10,89\,\% selon la pondération, contre 10,36\,\% avec les poids initiaux. Le franchissement d'un repère de 10\,\% dépend donc de la variante.

\textbf{Limites.} Les résultats sont conditionnels aux dates, à l'amplitude des variantes et aux catégories retenues. Les prix à collecte urbaine restent communs aux profils ruraux et urbains \cite{hcpIPC}. Les catégories partielles de l'élargi restent des approximations. Les quintiles sont définis par les dépenses au sein de chaque milieu et les montants en MAD reposent sur un budget illustratif identique. En l'absence de trajectoires de revenus et de quantités achetées, ces comparaisons portent sur le coût des postes couverts à composition fixe. Elles ne mesurent ni des comportements de substitution ni le pouvoir d'achat effectif.

\chapter*{Conclusion générale}
\addcontentsline{toc}{chapter}{Conclusion générale}
L'analyse combine deux sources officielles qui remplissent des rôles distincts. L'IPC établit la chronologie des variations de prix. L'ENNVM renseigne sur la répartition des dépenses des ménages selon leur milieu et leur quintile de dépenses. La période 2022--2023 concentre les pics de variation de l'IPC général dans l'échantillon. L'alimentation contribue à caractériser cet épisode par un glissement particulièrement élevé en février 2023.

Dans le modèle à poids fixes, la variation du coût et sa contribution au budget initial dépendent du périmètre. Pour l'ensemble national, de septembre 2022 à octobre 2026, la hausse historique est plus élevée dans le sous-panier prudent, tandis que la contribution au budget initial est plus forte dans le périmètre élargi. Le choc uniforme confirme que le taux du sous-panier est identique pour tous les profils et que le supplément dépend de la couverture. Le choc alimentaire montre que la charge marginale dépend de la part initiale du poste alimentaire. Ces résultats portent sur l'exposition des budgets sous les hypothèses choisies. Des revenus comparables, des quantités achetées et des prix locaux seraient nécessaires pour prolonger l'analyse du pouvoir d'achat et de la consommation.

L'introduction de ressources nominales hypothétiques précise le seuil de compensation. Une hausse de 5\,\% suffit exactement face au choc uniforme de 5\,\% dans les deux périmètres. Après le choc combiné, une hausse de 10\,\% suffit pour 14 profils dans le périmètre élargi et pour aucun dans le prudent. Ces résultats décrivent le nombre d'unités accessibles du sous-panier et dépendent des prix, du périmètre et des ressources supposées.

L'analyse de robustesse montre que le supplément alimentaire marginal reste supérieur pour Rural C1 à celui d'Urbain C5 sous les variantes de poids testées et sous l'hypothèse de prix communs. Le classement des périmètres pour la variation du coût historique change selon la fenêtre. Le dépassement de 10\,\% pour Rural C1 après le choc combiné dans l'élargi dépend aussi de la pondération. Ces constats circonscrivent les conclusions aux périodes, profils et variantes étudiés.

Les perspectives consistent à rapprocher plus finement les nomenclatures, à recueillir des prix ruraux comparables et à disposer de mesures répétées des revenus, dépenses et quantités achetées. Ces données permettraient de confronter les taux nominaux hypothétiques aux évolutions effectivement observées, de tester la stabilité des poids et d'étudier les substitutions ainsi que l'hétérogénéité au sein de chaque profil.

\appendix
\chapter{Correspondance détaillée des catégories}\label{ann:correspondance}
Les statuts « fiable », « partielle » et « absente » caractérisent la correspondance à l'échelle de la division mensuelle. Le statut « fiable » reprend la catégorie \texttt{proche} du CSV d'audit. L'exclusion du groupe 14 répond en plus à sa nature économique. Les prix associés aux groupes partiels servent uniquement au scénario élargi.
\begin{table}[!ht]
\centering\small
\setlength{\tabcolsep}{4pt}
\caption{Audit des correspondances ENNVM--IPC (groupes 01 à 07)}\label{tab:correspondance}
\begin{tabular}{@{}l>{\raggedright\arraybackslash}p{.235\textwidth}>{\raggedright\arraybackslash}p{.21\textwidth}l>{\raggedright\arraybackslash}p{.35\textwidth}@{}}
\hline
N° & Groupe ENNVM & Division IPC & Statut & Justification \\
\hline
01 & Aliments et boissons non alcoolisées & Produits alimentaires et boissons non alcoolisées & Fiable & DAMPG01 et IPC 011--012 : aliments (dont céréales, viandes et fruits) et boissons sans alcool. La division 01 est plus ciblée que l'agrégat Alimentation. \\
\hline
02 & Boissons alcoolisées, tabac et stupéfiants & Boissons alcoolisées, tabac et stupéfiants & Partielle & DAMPG02 publie une dépense de drogue (23), distincte de l'alcool et du tabac. Le détail IPC consulté montre 021 et 022, sans série identifiable pour les stupéfiants. \\
\hline
03 & Habillement et chaussures & Articles d'habillement et chaussures & Fiable & DAMPG03 et IPC 031--032 : vêtements, chaussures et services associés comme le nettoyage, la réparation ou la cordonnerie. \\
\hline
04 & Logement, eau, électricité, gaz et autres combustibles & Logement, eau, gaz, électricité et autres combustibles & Partielle & DAMPG04 comprend les loyers imputés (42). L'IPC documente les loyers payés par les locataires (041), avec recoupement sur l'eau, l'énergie et l'entretien. \\
\hline
05 & Ameublement, équipement ménager et entretien courant de la maison & Meubles, articles de ménage et entretien courant du foyer & Fiable & DAMPG05 et IPC 051--056 : ameublement, textiles, appareils, vaisselle, outils et entretien courant. \\
\hline
06 & Santé & Santé & Partielle & Médicaments, soins ambulatoires et hospitaliers se recoupent. DAMPG06 inclut aussi les groupes 64 et 65 (dont médecine traditionnelle), sans équivalent identifiable pour l'ensemble de ces dépenses dans le détail IPC consulté. \\
\hline
07 & Transport & Transports & Partielle & Véhicules, utilisation et transport de voyageurs se recoupent. DAMPG07 inclut aussi le fret et les services postaux (74), ces derniers figurant dans l'IPC 081. \\
\hline
\end{tabular}
\sourcetable{Fichier de correspondance du projet, groupes DAMPG de l'ENNVM et détail des classes de l'IPC.}
\end{table}
\clearpage
\begin{table}[!ht]
\centering\small
\setlength{\tabcolsep}{4pt}
\caption{Audit des correspondances ENNVM--IPC (groupes 08 à 14)}
\begin{tabular}{@{}l>{\raggedright\arraybackslash}p{.235\textwidth}>{\raggedright\arraybackslash}p{.21\textwidth}l>{\raggedright\arraybackslash}p{.35\textwidth}@{}}
\hline
N° & Groupe ENNVM & Division IPC & Statut & Justification \\
\hline
08 & Information et communication & Communications & Partielle & Recoupement sur la téléphonie. DAMPG08 comprend aussi le matériel informatique et audiovisuel que le détail IPC classe sous Loisirs et culture (091). \\
\hline
09 & Loisirs, sport et culture & Loisirs et culture & Partielle & Le sport, la culture et la papeterie se recoupent. L'IPC 091 inclut le matériel audiovisuel et informatique qui relève de DAMPG08 dans l'enquête. \\
\hline
10 & Enseignement & Enseignement & Fiable & DAMPG10 et IPC 101--105 : préscolaire et primaire, secondaire, postsecondaire non supérieur, supérieur et enseignement non défini par niveau. \\
\hline
11 & Services de restauration et d'hébergement (S) & Restaurants et hôtels & Partielle & Restauration, cantines et hébergement figurent des deux côtés. DAMPG11 détaille aussi les repas livrés et divers services dont la couverture précise reste incertaine. \\
\hline
12 & Assurances et services financiers & Aucune & Absente & Aucune série mensuelle autonome. Le classeur IPC annuel de détail situe les assurances transports et l'intermédiation financière dans Biens et services divers. \\
\hline
13 & Soins personnels, protection sociale et biens et services divers & Biens et services divers & Partielle & Recoupement sur les soins et effets personnels. DAMPG13 inclut la protection sociale et la division IPC recoupe aussi le groupe 12. \\
\hline
14 & Paiements fiscaux, transferts et remboursement des prêts (Non liés à l'exercice d'activité professionnelle) & Aucune & Absente & Poste distinct exclu des deux scénarios. Les paiements fiscaux, transferts et remboursements ne sont pas nécessairement des achats de biens et services. \\
\hline
\end{tabular}
\sourcetable{Fichier de correspondance du projet, groupes DAMPG de l'ENNVM et détail des classes de l'IPC.}
\end{table}


\chapter{Inventaire et contrôles de reproductibilité}\label{ann:trace}
\begin{description}\raggedright
\item[Données de prix.] Trois classeurs mensuels du HCP, base 100 en 2017, puis \path{ipc_maroc_mensuel.csv} (118 mois, 15 séries et une date). Le classeur annuel de détail des services divers documente certaines classes.
\item[Données de dépenses.] Classeur de tabulation ENNVM 2022--2023.
\par \path{depenses_menages_2022_2023.csv} (252 couples profil--poste) et \path{profils_budgetaires_2022_2023.csv} (18 profils).
\item[Correspondance.] Audit des 14 groupes dans \path{correspondance_depenses_ipc.csv}. Le notebook 03 contient le rapport de préparation, les couvertures et la définition de \texttt{simuler\_panier}.
\item[Analyse de l'IPC.] Le notebook 02 fournit les statistiques et quatre graphiques préfixés \texttt{02\_}. Le notebook 01 conserve la trame exploratoire, sans résultats exécutés.
\item[Simulations.] \path{04_cout_panier_historique.csv} contient 4\,248 lignes (118 mois, 18 profils, deux périmètres). \path{04_chocs_hypothetiques.csv} contient 72 lignes (deux chocs pour chaque couple profil--périmètre). \path{04_ponderations_sous_paniers.csv} contient 288 poids (18 profils, quatre postes prudents et douze élargis). Deux graphiques du notebook 04 complètent ces tableaux.
\item[Scénarios de ressources.] Le notebook 05 croise les coûts du notebook 04 avec quatre taux nominaux hypothétiques. Les fichiers \path{05_pouvoir_achat_historique.csv} (16\,992 lignes) et \path{05_pouvoir_achat_chocs.csv} (288 lignes) conservent la couverture et le seuil de maintien.
\item[Robustesse.] Le notebook 06 compare les dates de référence et les variantes de pondération alimentaire. Il enregistre ses tableaux et figures dans \path{outputs/robustesse/}.
\end{description}

Les contrôles du notebook 04 portent sur les clés, la complétude du calendrier, les valeurs finies, la conservation des poids sources, la normalisation à 100 en septembre 2022, et la concordance entre contribution directe et contribution obtenue via la couverture. Le choc uniforme produit exactement 5\,\% à la précision du calcul. Le supplément alimentaire marginal est identique entre périmètres pour un profil donné. Les taux historiques ont été confrontés à \texttt{simuler\_panier}. Les montants affichés dans le rapport sont arrondis après calcul.

\clearpage
\chapter{Tableaux complets des chocs hypothétiques}\label{ann:chocs}
Les colonnes de contribution donnent des points du budget initial. Les montants illustrent un même budget de 5\,000 MAD pour chaque profil. « Combiné » correspond à +5\,\% sur tous les postes couverts, puis +10\,\% supplémentaires sur le groupe 01.
\begin{table}[!htbp]
\centering\small
\caption{Contributions au budget initial, périmètre prudent}\label{tab:contributions-prudent}
\begin{tabular}{llrrr}
\hline
Milieu & Classe & Uniforme (pts) & Combiné (pts) & Surcroît aliment. (MAD) \\
\hline
National & C1 & 2{,}60 & 7{,}42 & 241{,}50 \\
National & C2 & 2{,}44 & 6{,}91 & 223{,}65 \\
National & C3 & 2{,}31 & 6{,}47 & 207{,}90 \\
National & C4 & 2{,}15 & 5{,}93 & 189{,}00 \\
National & C5 & 1{,}80 & 4{,}61 & 140{,}18 \\
National & Ensemble & 2{,}07 & 5{,}58 & 175{,}35 \\
Urbain & C1 & 2{,}25 & 6{,}36 & 205{,}28 \\
Urbain & C2 & 2{,}17 & 6{,}05 & 193{,}72 \\
Urbain & C3 & 2{,}12 & 5{,}81 & 184{,}80 \\
Urbain & C4 & 2{,}02 & 5{,}44 & 171{,}15 \\
Urbain & C5 & 1{,}73 & 4{,}31 & 129{,}15 \\
Urbain & Ensemble & 1{,}93 & 5{,}09 & 158{,}03 \\
Rural & C1 & 2{,}79 & 8{,}07 & 263{,}55 \\
Rural & C2 & 2{,}75 & 7{,}89 & 256{,}73 \\
Rural & C3 & 2{,}70 & 7{,}74 & 252{,}00 \\
Rural & C4 & 2{,}63 & 7{,}50 & 243{,}60 \\
Rural & C5 & 2{,}33 & 6{,}59 & 213{,}15 \\
Rural & Ensemble & 2{,}55 & 7{,}28 & 236{,}25 \\
\hline
\end{tabular}
\sourcetable{Simulations hypothétiques sur un budget initial de 5\,000 MAD.}
\end{table}

\begin{table}[!htbp]
\centering\small
\caption{Contributions au budget initial, périmètre élargi}\label{tab:contributions-elargi}
\begin{tabular}{llrrr}
\hline
Milieu & Classe & Uniforme (pts) & Combiné (pts) & Surcroît aliment. (MAD) \\
\hline
National & C1 & 4{,}90 & 9{,}73 & 241{,}50 \\
National & C2 & 4{,}86 & 9{,}33 & 223{,}65 \\
National & C3 & 4{,}79 & 8{,}95 & 207{,}90 \\
National & C4 & 4{,}72 & 8{,}51 & 189{,}00 \\
National & C5 & 4{,}50 & 7{,}31 & 140{,}18 \\
National & Ensemble & 4{,}66 & 8{,}17 & 175{,}35 \\
Urbain & C1 & 4{,}88 & 8{,}99 & 205{,}28 \\
Urbain & C2 & 4{,}81 & 8{,}68 & 193{,}72 \\
Urbain & C3 & 4{,}75 & 8{,}45 & 184{,}80 \\
Urbain & C4 & 4{,}67 & 8{,}09 & 171{,}15 \\
Urbain & C5 & 4{,}47 & 7{,}06 & 129{,}15 \\
Urbain & Ensemble & 4{,}62 & 7{,}78 & 158{,}03 \\
Rural & C1 & 4{,}92 & 10{,}19 & 263{,}55 \\
Rural & C2 & 4{,}89 & 10{,}02 & 256{,}73 \\
Rural & C3 & 4{,}85 & 9{,}89 & 252{,}00 \\
Rural & C4 & 4{,}79 & 9{,}66 & 243{,}60 \\
Rural & C5 & 4{,}63 & 8{,}90 & 213{,}15 \\
Rural & Ensemble & 4{,}76 & 9{,}48 & 236{,}25 \\
\hline
\end{tabular}
\sourcetable{Simulations hypothétiques sur un budget initial de 5\,000 MAD.}
\end{table}


\clearpage
\chapter{Pouvoir d'achat après le choc combiné}\label{ann:pa-chocs}
Le tableau de chaque périmètre présente les quatre hypothèses nominales et le seuil de maintien. Le groupe alimentaire subit une hausse totale de 15,5\,\% et les autres postes couverts de 5\,\%.
\begin{table}[!htbp]
\centering\small
\caption{Après le choc hypothétique combiné, périmètre prudent}\label{tab:pa-c-prudent}
\begin{tabular}{llrrrrrr}
\hline
Milieu & Classe & Couvert & Seuil & 0\,\% & +3\,\% & +5\,\% & +10\,\% \\
\hline
National & C1 & 51{,}9 & 14{,}31 & -12{,}52 & -9{,}89 & -8{,}14 & -3{,}77 \\
National & C2 & 48{,}8 & 14{,}17 & -12{,}41 & -9{,}78 & -8{,}03 & -3{,}65 \\
National & C3 & 46{,}2 & 14{,}00 & -12{,}28 & -9{,}65 & -7{,}89 & -3{,}51 \\
National & C4 & 43{,}0 & 13{,}79 & -12{,}12 & -9{,}48 & -7{,}73 & -3{,}33 \\
National & C5 & 36{,}1 & 12{,}77 & -11{,}32 & -8{,}66 & -6{,}89 & -2{,}45 \\
National & Ensemble & 41{,}4 & 13{,}47 & -11{,}87 & -9{,}23 & -7{,}47 & -3{,}06 \\
\noalign{\vskip 2pt}
Urbain & C1 & 45{,}0 & 14{,}12 & -12{,}38 & -9{,}75 & -7{,}99 & -3{,}61 \\
Urbain & C2 & 43{,}5 & 13{,}91 & -12{,}21 & -9{,}58 & -7{,}82 & -3{,}43 \\
Urbain & C3 & 42{,}3 & 13{,}74 & -12{,}08 & -9{,}44 & -7{,}68 & -3{,}29 \\
Urbain & C4 & 40{,}4 & 13{,}47 & -11{,}87 & -9{,}23 & -7{,}47 & -3{,}06 \\
Urbain & C5 & 34{,}6 & 12{,}47 & -11{,}08 & -8{,}42 & -6{,}64 & -2{,}19 \\
Urbain & Ensemble & 38{,}6 & 13{,}19 & -11{,}65 & -9{,}00 & -7{,}23 & -2{,}82 \\
\noalign{\vskip 2pt}
Rural & C1 & 55{,}9 & 14{,}43 & -12{,}61 & -9{,}99 & -8{,}24 & -3{,}87 \\
Rural & C2 & 55{,}1 & 14{,}32 & -12{,}53 & -9{,}90 & -8{,}15 & -3{,}78 \\
Rural & C3 & 54{,}0 & 14{,}33 & -12{,}54 & -9{,}91 & -8{,}16 & -3{,}79 \\
Rural & C4 & 52{,}6 & 14{,}26 & -12{,}48 & -9{,}86 & -8{,}11 & -3{,}73 \\
Rural & C5 & 46{,}5 & 14{,}17 & -12{,}41 & -9{,}78 & -8{,}03 & -3{,}65 \\
Rural & Ensemble & 51{,}0 & 14{,}26 & -12{,}48 & -9{,}86 & -8{,}11 & -3{,}73 \\
\hline
\end{tabular}
\sourcetable{Notebook 05. Couvert : part du budget initial en pourcentage. Seuil : croissance nominale de maintien en pourcentage. Les quatre dernières colonnes donnent la variation des unités accessibles (en \%).}
\end{table}

\begin{table}[!htbp]
\centering\small
\caption{Après le choc hypothétique combiné, périmètre élargi}\label{tab:pa-c-elargi}
\begin{tabular}{llrrrrrr}
\hline
Milieu & Classe & Couvert & Seuil & 0\,\% & +3\,\% & +5\,\% & +10\,\% \\
\hline
National & C1 & 98{,}0 & 9{,}93 & -9{,}03 & -6{,}30 & -4{,}48 & 0{,}06 \\
National & C2 & 97{,}1 & 9{,}61 & -8{,}76 & -6{,}03 & -4{,}20 & 0{,}36 \\
National & C3 & 95{,}9 & 9{,}34 & -8{,}54 & -5{,}79 & -3{,}97 & 0{,}61 \\
National & C4 & 94{,}5 & 9{,}00 & -8{,}26 & -5{,}50 & -3{,}67 & 0{,}92 \\
National & C5 & 90{,}1 & 8{,}11 & -7{,}50 & -4{,}73 & -2{,}88 & 1{,}75 \\
National & Ensemble & 93{,}2 & 8{,}76 & -8{,}06 & -5{,}30 & -3{,}46 & 1{,}14 \\
\noalign{\vskip 2pt}
Urbain & C1 & 97{,}7 & 9{,}20 & -8{,}43 & -5{,}68 & -3{,}85 & 0{,}73 \\
Urbain & C2 & 96{,}2 & 9{,}03 & -8{,}28 & -5{,}53 & -3{,}69 & 0{,}89 \\
Urbain & C3 & 95{,}1 & 8{,}89 & -8{,}16 & -5{,}41 & -3{,}57 & 1{,}02 \\
Urbain & C4 & 93{,}4 & 8{,}66 & -7{,}97 & -5{,}21 & -3{,}37 & 1{,}23 \\
Urbain & C5 & 89{,}5 & 7{,}89 & -7{,}31 & -4{,}53 & -2{,}68 & 1{,}96 \\
Urbain & Ensemble & 92{,}4 & 8{,}42 & -7{,}77 & -5{,}00 & -3{,}15 & 1{,}46 \\
\noalign{\vskip 2pt}
Rural & C1 & 98{,}3 & 10{,}36 & -9{,}39 & -6{,}67 & -4{,}86 & -0{,}33 \\
Rural & C2 & 97{,}8 & 10{,}25 & -9{,}30 & -6{,}58 & -4{,}76 & -0{,}23 \\
Rural & C3 & 97{,}0 & 10{,}20 & -9{,}25 & -6{,}53 & -4{,}72 & -0{,}18 \\
Rural & C4 & 95{,}8 & 10{,}09 & -9{,}16 & -6{,}44 & -4{,}62 & -0{,}08 \\
Rural & C5 & 92{,}7 & 9{,}60 & -8{,}76 & -6{,}02 & -4{,}20 & 0{,}37 \\
Rural & Ensemble & 95{,}2 & 9{,}96 & -9{,}06 & -6{,}33 & -4{,}51 & 0{,}03 \\
\hline
\end{tabular}
\sourcetable{Notebook 05. Couvert : part du budget initial en pourcentage. Seuil : croissance nominale de maintien en pourcentage. Les quatre dernières colonnes donnent la variation des unités accessibles (en \%).}
\end{table}


\clearpage
\phantomsection
\addcontentsline{toc}{chapter}{Références bibliographiques}
\renewcommand{\bibname}{Références bibliographiques}
\bibliographystyle{plain}
\bibliography{references}
\end{document}
